
\MakeThesisTitle
% Alternative fuer ein offizielles Titelblatt:
% \MakeOfficialTitlePage{offizielles-titelblatt.pdf}

\ifIncludeAbstract % J.3
  \addchap{Abstract}
  \textbf{Masterarbeiten} wird ein Abstract im Umfang von etwa ein, maximal zwei
  Seiten vorangestellt. Für alle anderen wissenschaftlichen Arbeiten ist dies
  nicht nötig! Ein Abstract stellt kurz und präzise die Kerninformationen zum
  Forschungshintergrund, zur Aufgabenstellung, zu den Leitthesen, der
  theoretischen Grundlage, der Methodik/Vorgehensweise, den Hauptergebnissen
  und dem Fazit der Arbeit (im Sinne eines Management Summary) dar. Es ist
  ergebnisorientiert zu verfassen und muss unabhängig vom Text verständlich
  sein. Die Leser sollen insbesondere die zentralen Erkenntnisse erfahren,
  die aus der Bearbeitung der Themenstellung gewonnen wurden.
\fi

\ThesisTableOfContents % J.4
\listoffigures % J.4
\listoftables % J.4

\addchap{\DEEN{Abkürzungsverzeichnis}{List of Abbreviations}} % J.5
Wenn Sie in Ihrer Arbeit Abkürzungen benutzen, sind diese in einem
Abkürzungsverzeichnis zu erläutern. Es wird nach den Kürzeln alphabetisch
aufsteigend sortiert. Allgemeinsprachliche Abkürzungen (wie vgl., z.\,B.,
u.\,a., usw., etc.) oder Abkürzungen, die üblicherweise in Quellenangaben
genutzt werden (wie ff., Nr., Jg., Hrsg., Aufl.) müssen \textit{nicht} im
Abkürzungsverzeichnis aufgenommen werden. Alle anderen Abkürzungen (z.\,B.
themen- und fachspezifische oder unternehmensinterne Abkürzungen, Namen von
Institutionen, Organisationen, Fachzeitschriften, Gesetzen oder andere
erklärungsbedürftige Abkürzungen) erfordern ausnahmslos einen Eintrag im
Abkürzungsverzeichnis.

Das Abkürzungsverzeichnis enthebt Sie nicht von der Aufgabe, eine Abkürzung
bei deren erstmaliger Verwendung im Text einzuführen: Bei der ersten Nennung
wird der Begriff ausgeschrieben und die Abkürzung in Klammern dahinter
gesetzt.

Im Folgenden eine beispielhafte Darstellung des
Abkürzungsverzeichnisses:

\begin{tabularx}{\textwidth}{@{}p{2.5cm}Y@{}}
  CP  & Credit Points \\
  HFH & Hamburger Fern-Hochschule \\
  PL  & Prüfungsleistung \\
  SSB & Studiengangsspezifische Bestimmungen
\end{tabularx}
